\section{Limits of finite tests on denominator powers}
\label{sec:denominator-obstructions}

The local obstruction behind the pole construction excludes odd negative
valuations.  It is natural to try to replace its many occurrences by
one product, or to impose finitely many such conditions directly on the
input.  The following results rule out two precise versions of those
strategies.  They do not identify the image of every pole equation with
the sets defined below; that identification would require an additional
converse theorem.

\begin{definition}
For an integer $m\ge2$, let
\[
 \mathcal D_m=\{q\in\mathbb Q_{>0}:
       \text{the positive reduced denominator of $q$ is an $m$th power}\}.
\]
Equivalently, $q\in\mathcal D_m$ if and only if $q>0$ and
$m\mid v_p(q)$ whenever $v_p(q)<0$.  In particular,
$\mathcal D_2$ consists of the positive rationals with square
denominator.
\end{definition}

\subsection{Multiplying local obstructions loses information}

\begin{proposition}\label{prop:denominator-monomial-rank}
Let $c_1,\ldots,c_r\in\mathcal D_2$, let
$A=(a_{ji})\in\operatorname{Mat}_{r\times n}(\mathbb Z)$, and put
\[
 M_j(b_1,\ldots,b_n)=c_j\prod_{i=1}^n b_i^{a_{ji}}.
\]
If $A$ has rank less than $n$ over $\mathbb F_2$, there are positive
rationals $b_1,\ldots,b_n$, some outside $\mathcal D_2$, such that
every $M_j(\mathbf b)$ belongs to $\mathcal D_2$.
Consequently an equivalence
\[
 \bigl(b_i\in\mathcal D_2\text{ for every }i\bigr)
 \quad\Longleftrightarrow\quad
 \bigl(M_j(\mathbf b)\in\mathcal D_2\text{ for every }j\bigr)
\]
on all positive rational inputs requires $r\ge n$ and full column
rank of $A$ modulo two.  These necessary conditions are not asserted
to be sufficient.
\end{proposition}
\begin{proof}
Choose a nonzero vector $\epsilon\in\ker(A:\mathbb F_2^n\to
\mathbb F_2^r)$ and a prime $p$ outside the numerator and denominator
supports of all $c_j$.  Set $b_i=p^{-\epsilon_i}$.  At least one
$b_i$ has reduced denominator $p$ and is not in $\mathcal D_2$.
But
\[
 v_p(M_j)=-\sum_i a_{ji}\epsilon_i
\]
is even.  Every other valuation of $M_j$ is the corresponding
valuation of $c_j$, so every negative valuation is even.  Positivity
is automatic.
\end{proof}

For example, the product of two inputs with denominator $p$ has square
denominator.  The proposition shows that choosing several fixed
monomial weights does not repair this cancellation unless the parity
conditions retain at least as much independent information as the
original conditions.

\subsection{Fixed rational pullbacks cannot characterize integrality}

A positive Boolean combination means a finite expression formed using
only conjunction and disjunction.  Rational functions are evaluated
only where their denominators are nonzero.  The result below permits
arbitrary fixed rational functions and different denominator exponents.

\begin{theorem}\label{thm:finite-denominator-tests}
Let $S\subseteq\mathbb Q$ be defined by a finite positive Boolean
combination of predicates of the following forms:
\[
 f(X)\in\mathcal D_m\quad(f\in\mathbb Q(T),\ m\ge2),
 \qquad g(X)=0,\qquad g(X)\ne0
 \quad(g\in\mathbb Q[T]).
\]
If $S$ contains all but finitely many rational integers, then there is
an integer $n\in S$ and infinitely many nonintegral elements of $S$
converging to $n$ in the real topology.  In particular no such finite
combination defines $\mathbb Z$ in $\mathbb Q$.
\end{theorem}

We first record the elementary prime-divisor fact needed in the proof.

\begin{lemma}\label{lem:simple-prime-values}
Let $h\in\mathbb Z[T]$ be nonconstant and squarefree over $\mathbb Q$.
Outside any prescribed finite set of primes, there is a prime $p$ and
an integer residue class $a\pmod{p^2}$ such that
\[
                        v_p(h(n))=1
              \quad\text{for every }n\equiv a\pmod{p^2}.
\]
One may simultaneously require finitely many polynomials coprime to
$h$ over $\mathbb Q$ to be nonzero modulo $p$ on this class.
\end{lemma}
\begin{proof}
A nonconstant integer polynomial has infinitely many prime divisors
of its nonzero integer values.  Here is a proof.  Translate its
argument so its constant term $c$ is nonzero.  If all such prime
divisors belonged to a finite set $\Pi$, choose $M$ divisible by
every prime in $\Pi$.  The nonconstant integer polynomial
\[
                         H(T)=h(cMT)/c
\]
has constant term one and satisfies $H(t)\equiv1\pmod M$.
Every prime divisor of a nonzero value of $H$ would belong to $\Pi$,
but none could divide that value.  Thus every nonzero value would
be $1$ or $-1$, a contradiction.

Now exclude the prescribed primes, the prime divisors of the leading
coefficient and discriminant of $h$, and those of the resultants with
the other polynomials.  A remaining prime divisor of a value supplies
a simple root $a_0$ modulo $p$.  Among its $p$ lifts modulo $p^2$,
exactly one is a root modulo $p^2$.  Any other lift has valuation one.
The resultant conditions give the last assertion.
\end{proof}

\begin{proof}[Proof of Theorem~\ref{thm:finite-denominator-tests}]
Write the defining expression in disjunctive normal form, as a finite
union of conjunctions, called branches in this proof.  A branch
containing a nonzero polynomial equality can contain only finitely
many integers; branches containing an identically false predicate
are empty.  Discard these branches and enlarge the finite exceptional
set of integers accordingly.  All remaining polynomial conditions
are identities or nonvanishing conditions.

Write each nonzero rational function in reduced form $P/Q$ over
$\mathbb Q[T]$.  Call a remaining branch bad if one of its predicates
$P(X)/Q(X)\in\mathcal D_m$ has an irreducible factor of $Q$ whose
multiplicity $e$ is not divisible by $m$.  For each bad branch choose
one such predicate and factor.  Exclude the prime supports of the
rational scalar factors and of all coefficient denominators in these
factorizations.  By
Lemma~\ref{lem:simple-prime-values}, choose a distinct prime $p$ and
a residue class modulo $p^2$ on which this factor has valuation one,
while $P$ and every other denominator factor are units.  Then
\[
                         v_p(P(n)/Q(n))=-e
\]
on that class, so the chosen predicate, and hence that branch, fails.
The Chinese remainder theorem combines these finitely many classes
into an infinite arithmetic progression on which all bad branches
fail.

Choose an integer $n$ in this progression outside the finite
exceptional set.  Since $n\in S$, it belongs to a remaining branch
which is not bad.  Every denominator factor in each predicate of
this branch has multiplicity divisible by its exponent $m_j$.
After absorbing a nonzero rational constant into the numerator,
write all its denominator predicates as
\[
              f_j(X)=P_j(X)/B_j(X)^{m_j}\in\mathcal D_{m_j},
              \qquad\gcd(P_j,B_j)=1.
\]
These rational functions are defined and positive at $n$.

Let $\Pi$ be a finite set of primes containing all coefficient
denominators, leading-coefficient supports, nonzero resultant
supports, and the numerator and denominator supports of the
nonzero rational numbers $f_j(n)$.  Enlarge it if needed so that
outside $\Pi$ all the displayed polynomials are integral with
unit leading coefficient and each pair $(P_j,B_j)$ has unit
resultant.  For each $q\in\Pi$, continuity in $\mathbb Q_q$
gives an integer $e_q\ge1$ such that
\[
 x-n\in q^{e_q}\mathbb Z_q
       \quad\Longrightarrow\quad
 v_q(f_j(x))=v_q(f_j(n))\quad\text{for every }j.
\]
Choose a positive integer $M$ divisible by every $q^{e_q}$,
and put $L=\operatorname{lcm}(m_1,\ldots,m_s)$.  If there are
no denominator predicates, set $L=2$.  For any sufficiently large
prime $p\notin\Pi$ not dividing $M$, define
\[
                              x_p=n+M/p^L.
\]
These are distinct nonintegers and converge to $n$ in $\mathbb R$.

We check every denominator condition of the chosen branch.  At
$q\in\Pi$ the valuation agrees with its value at $n$, which
satisfies that condition.  At $p$ we have $v_p(x_p)=-L$; the unit
leading coefficients give
\[
 v_p(f_j(x_p))=-L\deg P_j+m_jL\deg B_j,
\]
a multiple of $m_j$.  At a prime $q\notin\Pi\cup\{p\}$,
the number $x_p$ is integral.  If $B_j(x_p)$ is a unit, then
$v_q(f_j(x_p))\ge0$.  Otherwise the unit resultant makes
$P_j(x_p)$ a unit and
\[
                      v_q(f_j(x_p))=-m_jv_q(B_j(x_p)),
\]
again a permitted valuation.  Real continuity makes all $f_j(x_p)$
positive for sufficiently large $p$, and preserves the finitely
many polynomial nonvanishing conditions.  Thus all sufficiently
large choices of $p$ give $x_p\in S\setminus\mathbb Z$.
\end{proof}

\begin{corollary}\label{cor:boolean-denominator-tests}
The conclusion of Theorem~\ref{thm:finite-denominator-tests} also
holds for arbitrary finite Boolean combinations of the listed
predicates, allowing negation.
\end{corollary}
\begin{proof}
Use disjunctive normal form with both positive and negative denominator
literals.  Negated polynomial equalities and inequalities are already
among the allowed predicates.  Discard finite branches as before, and
use the same Chinese remainder construction to eliminate every branch
with a bad positive denominator literal.  Choose the resulting
admitted integer $n$ also away from the finitely many zeros and poles
of all nonzero rational functions in the formula.

Consider a negative literal $f(n)\notin\mathcal D_m$ in the chosen
branch.  An identically zero $f$ makes it automatically true.  Otherwise
either $f(n)<0$, which persists in a real neighborhood, or $f(n)>0$
and there is a prime $q$ with $v_q(f(n))<0$ not divisible by $m$.
The latter failure persists in a sufficiently small $q$-adic
neighborhood of $n$.  Add one such prime for each relevant negative
literal to $\Pi$, and choose $M$ sufficiently divisible to preserve
all those valuations as well.  The same numbers $x_p=n+M/p^L$
then satisfy every positive and negative literal for all sufficiently
large permitted primes $p$.  They are nonintegral and converge to $n$.
\end{proof}

The theorem is deliberately about predicates on fixed rational
functions of the input, with no extra existential coordinates.
It does not rule out auxiliary elliptic points, variable rational
functions, or other arithmetic predicates.  It also does not replace
the independent proof needed to identify the exact image of a given
pole system.  Its conclusion is that even arbitrarily complicated
fixed rational pullbacks of perfect-denominator conditions cannot
alone remove every nonintegral rational number.
